% Canonical atom manuscript integration fragment.
% === STRUCTURAL BLOCK ===
\begin{theorem}[Canonical atom normal form]
\label{theorem-canonical-atom-normal-form}
Let $F$ be a finite triple system without isolated points satisfying the three
intrinsic conditions in Theorem~A.  Its hyperedges admit a canonical partition
into \emph{atoms} with the following properties.
\begin{enumerate}
\item Every atom is either one triple or $J^+$ for a finite $2$-connected
simple bipartite graph $J$.
\item Two distinct atoms meet in at most one point.  The bipartite incidence
graph between atoms and points belonging to at least two atoms is a forest.
\item The atom partition, the nontrivial cores $J$ up to graph isomorphism, and
the atom--point forest are determined by $F$.
\end{enumerate}
Conversely, every forest assembly of these atoms by disjoint unions and
one-point amalgamations satisfies the intrinsic conditions.
\end{theorem}
\begin{proof}
Let $L=I(F)$.  A \emph{cyclic block} is a maximal $2$-connected subgraph of
$L$ containing a cycle.  Since every hyperedge-node has degree three and is
incident with a bridge, its degree inside a cyclic block is two.  Suppressing
the hyperedge-nodes in such a block $C$ gives a finite graph $J_C$ on the
point-nodes of $C$.  It is simple by linearity and $2$-connected by standard
block theory.  Its cycles correspond to Berge cycles of the same lengths, so
$J_C$ is bipartite.

For a hyperedge-node $e\in C$, the third incidence is a bridge.  Its point
endpoint lies outside $C$, since otherwise that incidence would lie on a
cycle.  These third points are pairwise distinct as $e$ varies in $C$;
otherwise a path inside $C$ between two hyperedge-nodes, completed through the
common third point, would put both bridge incidences on a cycle.  The
hyperedges represented in $C$ therefore form exactly $J_C^+$.

A hyperedge-node outside every cyclic block has no nonbridge incidence and
contributes one single-triple atom.  Thus the atoms partition the hyperedges.
Two atoms cannot share two points, since paths inside the two atoms would
produce a Levi cycle crossing two cyclic blocks.  Similarly, a cycle in the
atom--shared-point incidence graph expands to a closed Levi walk and hence to
a Levi cycle crossing atom boundaries.  The incidence graph is therefore a
forest.  Cyclic blocks and the block--cut forest are canonical, proving the
forward statement.  The converse follows because an atomic expansion is
linear, has a private bridge at each hyperedge-node, and has only even Berge
cycles, while a forest of one-point sums preserves these properties.
\end{proof}

\begin{corollary}[Minimal generators and one-point indecomposables]
\label{corollary-minimal-generators-indecomposables}
The class $\mathcal B$ is generated, up to isomorphism, by finite edgeless
systems, one triple, and $J^+$ for finite $2$-connected bipartite graphs $J$,
under finite disjoint unions and one-point amalgamations.  Moreover, a
connected reduced obligatory system with at least one hyperedge is
one-point indecomposable if and only if it is one triple or $J^+$ for a finite
$2$-connected bipartite graph $J$.
\end{corollary}
\begin{proof}
The theorem decomposes every reduced member of $\mathcal B$ into the displayed
atoms; isolated points are restored by an edgeless factor.  If the atom forest
has at least two atom-nodes, a leaf atom separates from the rest at its unique
shared point.  Conversely, one triple is indecomposable, and if $J$ is
$2$-connected then deleting any point of $J^+$ cannot separate its
hyperedge-nodes into two nonempty classes: a private point is a leaf, while a
core point leaves $J$ connected.  Hence $J^+$ is indecomposable.
\end{proof}
% === FINITE SPECTRA BLOCK ===
\subsection{Canonical atoms and structural spectra}
\label{canonical-atoms-and-structural-spectra}

For a connected reduced obligatory system $F$, put
\[
s=|V(F)|-|E(F)|,
\qquad
\beta=2|E(F)|-|V(F)|+1.
\]
Thus $s$ is the order of a connected bipartite shadow and $\beta$ is its
cycle rank.  Let $a(F)$ denote the number of canonical atoms from
Theorem~\ref{theorem-canonical-atom-normal-form}.

\begin{lemma}[Minimum order at fixed cyclic rank]
\label{lemma-minimum-order-fixed-cyclic-rank}
Let $J$ be a finite $2$-connected simple bipartite graph of cycle rank
$r=|E(J)|-|V(J)|+1\ge1$.  Then
\[
|V(J)|\ge 2+q(r).
\tag{10.12}
\]
For $r\ge2$, every order $v\ge2+q(r)$ occurs.  For $r=1$, the possible orders
are exactly the even integers $v\ge4$.
\end{lemma}
\begin{proof}
Writing $v=|V(J)|$, the bipartite extremal bound gives
\[
r\le\left\lfloor\frac{v^2}{4}\right\rfloor-v+1
 =\left\lfloor\frac{(v-2)^2}{4}\right\rfloor,
\]
which is equivalent to (10.12).  If $v$ is even, start with the spanning cycle
$C_v$ in the balanced complete bipartite graph and add edges until the desired
rank is reached.  If $v$ is odd, start with $C_{v-1}$ and add the remaining
vertex with two neighbours in the opposite bipartition class; this is
$2$-connected and has rank two.  Add arbitrary missing cross-edges.  These
constructions realise every allowed $r\ge2$.  A $2$-connected graph of rank one
is a cycle, which is bipartite exactly when its order is even.
\end{proof}

\begin{theorem}[Exact canonical atom-count spectrum]
\label{theorem-exact-canonical-atom-count-spectrum}
Let $F$ range over connected reduced obligatory systems with fixed shadow
order $s$ and cycle rank $\beta$.  The possible values $k=a(F)$ are exactly
\[
\begin{array}{ll}
\beta=0:
  & k=s-1,\\[1mm]
\beta=1:
  & 1\le k\le s-3,\quad k\equiv s+1\pmod 2,\\[1mm]
\beta\ge2:
  & 1\le k\le s-1-q(\beta).
\end{array}
\tag{10.13}
\]
For $\beta\ge1$, the maximum is
\[
a_{\max}(s,\beta)=s-1-q(\beta).
\tag{10.14}
\]
Every maximiser has one minimum-order cyclic atom carrying the full rank
$\beta$, and every other atom is one triple.  The parity obstruction in the
rank-one line is the only gap in any atom-count spectrum.
\end{theorem}
\begin{proof}
Let the atom cores be $J_i$, with $v_i=|V(J_i)|$,
$e_i=|E(J_i)|$, and $r_i=e_i-v_i+1$.  Since a tree of $k$ atoms makes
$k-1$ one-point identifications,
\[
s=1+\sum_i(v_i-1),
\qquad
\beta=\sum_i r_i.
\tag{10.15}
\]
If $\beta=0$, every atom is one triple and (10.15) gives $k=s-1$.
If $\beta=1$, there is one cyclic atom, whose core is an even cycle of order
$v\ge4$; the remaining atoms are single triples, so
$s=k+v-1$.  This is exactly the middle line of (10.13).

Suppose $\beta\ge2$.  Lemma~\ref{lemma-minimum-order-fixed-cyclic-rank} and
\[
q(r_1)+q(r_2)\ge q(r_1+r_2)+1
\]
give
\[
s\ge1+k+q(\beta),
\]
so $k\le s-1-q(\beta)$.  Conversely, for every integer in the displayed
interval choose one cyclic core of rank $\beta$ and order $s-k+1$, supplied by
the lemma, and attach $k-1$ single-triple atoms.  Equality in (10.14) forces
one positive-rank atom and minimum order, proving the rigidity statement.
\end{proof}

For $s\ge3$, define
\[
\alpha(s)=
\begin{cases}
s,&s\text{ even},\\
s+1,&s\text{ odd},
\end{cases}
\qquad
\delta(s)=\left\lfloor\frac{(s-1)^2}{4}\right\rfloor+1.
\]

\begin{theorem}[Indecomposability phase diagram]
\label{theorem-indecomposability-phase-diagram}
Apart from the single triple $(|E(F)|,|V(F)|)=(1,3)$, a connected reduced
one-point-indecomposable obligatory system with shadow order $s$ and $m$
hyperedges exists if and only if
\[
s\ge4,
\qquad
\alpha(s)\le m\le\left\lfloor\frac{s^2}{4}\right\rfloor.
\tag{10.16}
\]
A connected reduced one-point-decomposable obligatory system exists if and
only if
\[
s\ge3,
\qquad
s-1\le m\le\delta(s).
\tag{10.17}
\]
Consequently the feasible connected parameters split into three exact zones:
\[
\begin{array}{rcl}
s-1\le m<\alpha(s)
 &\Longrightarrow& \text{every system is decomposable},\\
\alpha(s)\le m\le\delta(s)
 &\Longrightarrow& \text{both behaviours occur},\\
\delta(s)<m\le\lfloor s^2/4\rfloor
 &\Longrightarrow& \text{every system is indecomposable}.
\end{array}
\tag{10.18}
\]
For $m\ge2$, the last zone is equivalently
\[
s\le\left\lceil2\sqrt{m-1}\right\rceil.
\tag{10.19}
\]
\end{theorem}
\begin{proof}
By Corollary~\ref{corollary-minimal-generators-indecomposables}, the nontrivial
indecomposable systems are $J^+$ with $J$ finite, simple, bipartite and
$2$-connected.  If $|V(J)|=s$ and $|E(J)|=m$, minimum degree two gives
$m\ge s$; equality makes $J$ a cycle and therefore forces $s$ even.  The
upper bound is the usual $m\le\lfloor s^2/4\rfloor$.  Even cycles, and for odd
$s$ a cycle on $s-1$ vertices plus one new vertex with two neighbours, give
the lower endpoints; adding cross-edges fills the intervals.

A decomposable system admits a connected bipartite shadow with a cut vertex.
If its bipartition sizes are $a+b=s$ and the cut vertex lies in the $a$-side,
then deleting it gives components with sizes $(a_i,b_i)$ and
\[
m\le b+\sum_i a_i b_i
 \le b+(a-1)(b-1)=ab-a+1.
\]
The symmetric estimate gives $m\le ab-\min(a,b)+1\le\delta(s)$.  The bound is
sharp: take a connected bipartite graph on $s-1$ vertices with $m-1$ edges and
adjoin one leaf.  This also fills every value in (10.17).  The three zones and
the elementary inversion (10.19) follow.
\end{proof}

\begin{corollary}[Rigidity on the structural boundaries]
\label{corollary-structural-boundary-rigidity}
Let $F$ be connected, reduced and obligatory, with shadow order $s$.
\begin{enumerate}
\item If $F$ is indecomposable, $s$ is even and $m=s$, then
$F\cong C_s^+$.
\item If $F$ is indecomposable, $s$ is odd and $m=s+1$, then
$F\cong\Theta(2a,2b,2c)^+$ for positive $a\le b\le c$ with
$a+b+c=(s+1)/2$.
\item If $F$ is indecomposable and $m=\lfloor s^2/4\rfloor$, then
$F\cong K_{\lfloor s/2\rfloor,\lceil s/2\rceil}^+$.
\item If $F$ is decomposable and $m=\delta(s)$, then $F$ is a one-point
amalgamation of $K_{a,b}^+$ and one triple, where
$a+b=s-1$ and $|a-b|\le1$.
\end{enumerate}
Here $\Theta(p,q,r)$ is the theta graph with three internally disjoint paths
of lengths $p,q,r$ and common endpoints.
\end{corollary}
\begin{proof}
The first and third assertions are the equality cases of minimum degree and
the bipartite extremal bound.  In the second case
$\sum_v(\deg v-2)=2$.  A single degree-four vertex would be a cut vertex, so
there are exactly two degree-three vertices and all others have degree two;
the three resulting paths have the same parity, hence are all even.  Equality
in the cut-vertex estimate used above forces a balanced complete bipartite
graph on $s-1$ vertices plus one leaf, which translates to the fourth
assertion.
\end{proof}

\begin{proposition}[Componentwise atom-count spectrum]
\label{proposition-componentwise-atom-count-spectrum}
Let $F$ be reduced and obligatory with $c$ connected components, shadow order
$s=|V(F)|-|E(F)|$, cycle rank
$\beta=2|E(F)|-|V(F)|+c$, and $k$ canonical atoms.  For every feasible
$(s,\beta,c)$, the possible values of $k$ are exactly
\[
\begin{array}{ll}
\beta=0:
  & k=s-c,\\[1mm]
\beta=1:
  & c\le k\le s-c-2,\quad k\equiv s-c\pmod2,\\[1mm]
\beta\ge2:
  & c\le k\le s-c-q(\beta).
\end{array}
\tag{10.20}
\]
The parameters themselves are feasible exactly when $s\ge2c$ for $\beta=0$
and $s\ge2c+q(\beta)$ for $\beta\ge1$.
\end{proposition}
\begin{proof}
If the atom cores are $J_i$, a forest with $k$ atom-nodes and $c$ components
makes $k-c$ identifications.  Hence
\[
s=c+\sum_i(|V(J_i)|-1),
\qquad
\beta=\sum_i\bigl(|E(J_i)|-|V(J_i)|+1\bigr).
\]
The connected proof applies componentwise.  Necessity gives the lower atom
floor $k\ge c$ and the three upper formulas.  For sufficiency, take $c-1$
single-triple components and realise the corresponding connected atom count
in the remaining component.  The feasibility thresholds follow by
concentrating all positive rank in one minimum-order cyclic atom.
\end{proof}
